\section{Results and Evaluation}
\label{sec:results}

Everything below is the untouched test window --- orders placed 30 June to 30
August 2018 --- scored once, for models already chosen. Cross-validated and validation columns sit beside
it because the disagreement between them is part of the result;
Appendix~\ref{app:tvt} adds the in-sample training fit of the deployed models,
and Section~\ref{sec:results-segments} carries the error analysis by month and
region. Figures are in Appendix~\ref{app:figures}.

\subsection{Problem 1A --- lead time}
\label{sec:results-reg}

\begin{table}[htbp]
\centering
\caption{Lead-time regression: two reference points that are not models, then
each family's complexity ladder. CV MAE is the mean over the five expanding
folds with its fold-to-fold standard deviation. \best{Bold} is the selected
and deployed model. Figure~\ref{fig:ladder}.}
\label{tab:reg}
\small
\setlength{\tabcolsep}{3.5pt}
\resizebox{\linewidth}{!}{%
\begin{tabular}{@{}llrrrrr@{}}
\toprule
\textbf{Model} & \textbf{Family} & \textbf{CV MAE $\pm$ sd} & \textbf{Val MAE} &
\textbf{Test MAE} & \textbf{RMSE} & \textbf{$R^2$} \\
\midrule
Mean of the training period      & baseline & ---  & ---  & 6.645 & 7.572 & $-0.990$ \\
Olist's own delivery promise     & baseline & ---  & ---  & 9.788 & 12.155 & $-4.130$ \\
\midrule
LinearRegression \emph{(rung 0)} & A linear & 5.242 $\pm$ 1.00 & 4.965 & 4.181 & 5.639 & $-0.104$ \\
Ridge \emph{(rung 1)}            & A linear & 5.201 $\pm$ 0.95 & 5.164 & 4.309 & 5.775 & $-0.158$ \\
DecisionTree \emph{(rung 0)}     & B tree   & 5.539 $\pm$ 0.91 & 5.927 & 6.116 & 8.664 & $-1.606$ \\
RandomForest \emph{(rung 1)}     & B tree   & 5.457 $\pm$ 0.90 & 5.513 & 4.956 & 6.169 & $-0.321$ \\
HistGradientBoosting \emph{(rung 2)} & B tree & 5.495 $\pm$ 0.81 & 4.802 & 4.141 & 5.517 & $-0.057$ \\
\best{VotingRegressor (A+B)}     & C ensemble & 5.315 $\pm$ 0.87 & \best{4.784} & \best{3.898} & \best{5.342} & \best{$+0.009$} \\
StackingRegressor (A+B)          & C ensemble & 5.443 $\pm$ 0.84 & 4.792 & 4.121 & 5.430 & $-0.024$ \\
\bottomrule
\end{tabular}}
\end{table}

The selected model puts mean absolute error at \textbf{3.90 days}: \textbf{41.3\%
below} a constant fitted on the training period and \textbf{60.2\% below the
platform's own promised date}. That second comparison needs care. A promised
date is a commercial decision, padded because a broken promise costs returns and
retention \citep{cui2024promise} --- an excellent feature and a poor forecast ---
and the forecast is not a promise either: on its own it is exceeded by 31.9\%
of delivered test orders, which is why Section~\ref{sec:stakeholders} adds a
margin.

Three things matter more than the winner. \textbf{Regularisation does not pay:}
Ridge is \emph{worse} than plain least squares on validation (5.164 against
4.965) and test (4.309 against 4.181). \textbf{Inside family B complexity does
pay}, monotonically on validation: 5.927 $\rightarrow$ 5.513 $\rightarrow$ 4.802.
And \textbf{only the ensemble has a positive $R^2$}; the others' negative values
are the signature of a level shift rather than of useless models, and the reason
MAE leads \citep{hyndman2006accuracy}. The blend beats both members (3.898
against 4.181 and 4.141) because a tree cannot predict outside the range of its
training targets, test-window lead times are shorter than anything common in
training, and a linear model extrapolates the level (Figure~\ref{fig:extrap});
the 1:2 weights are fixed, not tuned.

\subsection{Problem 1B --- late against the promise}
\label{sec:results-late}

\begin{table}[htbp]
\centering
\caption{The late-order flag, best of each route; the full seven-candidate table
is in \code{report/results/clf\_p1b.csv}. No-skill PR-AUC floor 0.074.
Figure~\ref{fig:routes}.}
\label{tab:late}
\footnotesize
{\setlength{\tabcolsep}{3pt}
\resizebox{\linewidth}{!}{%
\begin{tabular}{@{}llrrrrrrr@{}}
\toprule
\textbf{Route} & \textbf{Model} & \textbf{ROC-AUC} & \textbf{PR-AUC} &
\textbf{lift} & \textbf{bal.\ acc} & \textbf{prec.} & \textbf{recall} & \textbf{F1} \\
\midrule
direct  & best of five (HistGB) & 0.584 & 0.092 & 1.24$\times$ & 0.524 & 0.088 & 0.279 & 0.134 \\
derived & predicted $>$ promised, unfitted & 0.569 & 0.109 & 1.47$\times$ & 0.569 & 0.274 & 0.175 & 0.214 \\
\best{derived} & \best{regression margin, calibrated} & \best{0.782} & \best{0.214} & \best{2.89$\times$} & \best{0.705} & 0.137 & \best{0.826} & \best{0.235} \\
\bottomrule
\end{tabular}}}
\end{table}

The deliverable is the derived route, and that follows from the framing fixed in
Section~\ref{sec:problem} rather than from the numbers: Problem 1 is the brief's
dual framing, so the flag \emph{is} the regression's own answer compared against
the promised date. The five direct classifiers answer a different question ---
does the flag need a model of its own? It does not. The derived route reaches
\textbf{2.89$\times$} the no-skill floor where the best direct classifier reaches
1.24$\times$, and it costs no additional model: the margin
$\hat{y} - \text{promised}$ is mapped to a probability by a two-parameter
logistic \citep{platt1999probabilistic} fitted on validation. Even the
zero-parameter rule --- flag whenever the model expects a miss --- beats every
fitted classifier on PR-AUC.

\paragraph{The selection protocol got this one wrong.} Applied to these seven
candidates the rolling-origin rule of Section~\ref{sec:cv} ranks a direct
classifier first: on the in-training windows the candidates sit within noise of
one another, while on the forward window the derived route leads by a factor of
2.3. Those windows are dominated by the February--March 2018 slowdown and the
strike (Section~\ref{sec:eda-lateness}); a classifier can fit a regime, the
derived route encodes a mechanism (Appendix~\ref{app:protocol}). We report the
disagreement rather than credit the rule with this choice.

\subsection{Problem 2 --- low-review risk}
\label{sec:results-review}

\begin{table}[htbp]
\centering
\caption{Low-review classification at delivery. CV is the mean PR-AUC over the
five expanding folds with its fold-to-fold standard deviation
(Figure~\ref{fig:cvfolds}; every fold in Table~\ref{tab:cvfolds}). Accuracy is
reported in the text; the majority-class rule scores 0.890, flagging no one.
No-skill PR-AUC floor 0.110. \best{Bold} is the selected and deployed model.
Figure~\ref{fig:lowreview}.}
\label{tab:review}
\footnotesize
{\setlength{\tabcolsep}{2.2pt}
\resizebox{\linewidth}{!}{%
\begin{tabular}{@{}llrrrrrrrrrr@{}}
\toprule
\textbf{Model} & \textbf{Fam.} & \textbf{CV $\pm$ sd} & \textbf{Val} & \textbf{ROC} &
\textbf{PR-AUC} & \textbf{lift} & \textbf{bal.\ acc} & \textbf{prec.} &
\textbf{recall} & \textbf{F1} & \textbf{Brier} \\
\midrule
LogisticRegression \emph{(r0)}   & A & 0.465 $\pm$ .095 & 0.320 & 0.659 & 0.291 & 2.63$\times$ & 0.636 & 0.328 & 0.365 & 0.345 & 0.090 \\
LogisticRegression ($\ell_2$) \emph{(r1)} & A & 0.466 $\pm$ .095 & 0.320 & 0.659 & 0.291 & 2.63$\times$ & 0.636 & 0.317 & 0.371 & 0.342 & 0.090 \\
DecisionTree \emph{(r0)}         & B & 0.543 $\pm$ .092 & 0.389 & 0.722 & 0.396 & 3.59$\times$ & 0.677 & 0.562 & 0.391 & 0.461 & 0.080 \\
RandomForest \emph{(r1)}         & B & 0.568 $\pm$ .089 & 0.403 & 0.744 & 0.405 & 3.67$\times$ & 0.680 & 0.576 & 0.396 & 0.470 & 0.078 \\
HistGradientBoosting \emph{(r2)} & B & 0.568 $\pm$ .088 & 0.407 & 0.749 & 0.415 & 3.76$\times$ & 0.680 & 0.579 & 0.396 & 0.471 & 0.078 \\
VotingClassifier (A+B)           & C & 0.561 $\pm$ .084 & 0.404 & 0.747 & 0.407 & 3.69$\times$ & 0.683 & 0.477 & 0.422 & 0.448 & 0.079 \\
\best{StackingClassifier (A+B)}  & C & 0.568 $\pm$ .087 & \best{0.409} & 0.749 & 0.414 & 3.75$\times$ & 0.681 & 0.575 & 0.398 & 0.470 & 0.078 \\
\bottomrule
\end{tabular}}}
\end{table}

The selected model reaches \textbf{ROC-AUC 0.749}, \textbf{balanced accuracy
0.681} and \textbf{PR-AUC 0.414 --- 3.75$\times$ the no-skill floor}. Its raw
accuracy is 0.901 against a majority-class rule scoring 0.890, which is exactly
why accuracy is printed but never led with. The tree family beats the linear one
by 1.43$\times$ on PR-AUC --- but a \emph{single depth-limited decision tree}
already reaches 0.396 of the eventual 0.414, so 96\% of the achievable signal
comes from one interpretable tree. And the model does not memorise buyers: on the
12{,}617 test orders from customers never seen in training it scores PR-AUC
\textbf{0.415} against 0.414 overall. On the 12{,}571 delivered parcels alone it
scores ROC-AUC 0.711 and PR-AUC 0.333, 3.5 times their 9.6\% floor; the 284
reviewed orders that never arrived, 73\% of them rated 1--2 stars, raise the
overall figures (Appendix~\ref{app:timing}).

\paragraph{What the fold-to-fold variation says.}
The Phase 2 evaluation asked for the spread behind each cross-validated number;
Figure~\ref{fig:cvfolds} and Table~\ref{tab:cvfolds} show every candidate on
every window. The spread, a standard deviation of about 0.09, is almost entirely
a property of the \emph{windows}: the mean over candidates runs from 0.376 on
the first fold to 0.630 on the fourth, while within a window the candidates' sd has
a median of 0.045. So the forest, the boosted model and the stack (CV means of
0.568 apiece) cannot be separated by cross-validation, the case for the
rank-based rule of Section~\ref{sec:cv}; the linear family is behind the trees on
every window, so that conclusion is robust to the spread.

\paragraph{Class weighting: the ablation Phase 1 promised.} On the best tree
model, weighting the loss leaves ranking unchanged (ROC-AUC 0.749 to 0.747,
PR-AUC 0.415 to 0.411) while the Brier score degrades by 71\% (0.0778 to
0.1331) and mean predicted risk inflates from 0.119 to \textbf{0.337} against an
actual 0.110 (Figure~\ref{fig:calibration}). We therefore fit on the observed
prior and move the threshold: tuned on validation by maximising F1, it lands at
0.19 rather than 0.50 and raises test F1 from 0.414 to \textbf{0.470};
calibration holds where the volume is (a gap under 0.013 in the bins holding
11{,}892 of the 12{,}855 orders).

\subsection{Precision@$K$ and Recall@$K$ at a fixed capacity}
\label{sec:results-capacity}

\begin{table}[htbp]
\centering
\caption{What a ranked list delivers at a fixed capacity, on the test window.
$K$ is the share of reviewed (late flag: delivered) orders acted on; the daily rows group them by
purchase day, as on the slides, while the list a team would work groups
delivered, not-yet-rated parcels by delivery day (text). Break-even precision is
$c_{\text{action}}/c_{\text{miss}}$ (Section~\ref{sec:cost}), if every action works.
Figure~\ref{fig:capacity}.}
\label{tab:capacity}
\footnotesize
{\setlength{\tabcolsep}{2.6pt}
\resizebox{\linewidth}{!}{%
\begin{tabular}{@{}lrrrrrr@{}}
\toprule
 & \multicolumn{3}{c}{\textbf{Review risk, at delivery}} & \multicolumn{3}{c}{\textbf{Late flag, at checkout: derived / direct}} \\
\cmidrule(lr){2-4}\cmidrule(lr){5-7}
\textbf{Capacity $K$} & \textbf{precision} & \textbf{recall} & \textbf{lift} & \textbf{precision} & \textbf{recall} & \textbf{lift} \\
\midrule
top 1\%  & 0.705 & 0.064 & 6.4$\times$ & 0.270 / 0.111 & 0.036 / 0.015 & 3.6 / 1.5$\times$ \\
top 5\%  & 0.678 & 0.307 & 6.1$\times$ & 0.274 / 0.104 & 0.185 / 0.071 & 3.7 / 1.4$\times$ \\
top 10\% & 0.475 & 0.431 & 4.3$\times$ & 0.258 / 0.093 & 0.348 / 0.126 & 3.5 / 1.3$\times$ \\
top 20\% & 0.300 & 0.544 & 2.7$\times$ & 0.211 / 0.089 & 0.568 / 0.241 & 2.8 / 1.2$\times$ \\
top 30\% & 0.236 & 0.641 & 2.1$\times$ & 0.175 / 0.086 & 0.707 / 0.349 & 2.4 / 1.2$\times$ \\
\midrule
10 a day (4.8\%) & 0.572 & 0.248 & 5.2$\times$ & --- & --- & --- \\
20 a day (9.5\%) & 0.440 & 0.378 & 4.0$\times$ & --- & --- & --- \\
\midrule
break-even precision & \multicolumn{3}{c}{0.125 \ (R\$5 : R\$40)} & \multicolumn{3}{c}{0.300 \ (R\$12 : R\$40)} \\
\bottomrule
\end{tabular}}}
\end{table}

Table~\ref{tab:capacity} reads the predictions of Tables~\ref{tab:late} and
\ref{tab:review} as a team with limited capacity would. \textbf{Capacity, not
the threshold, is the binding constraint for the review flag.} At a capacity of 5\% of reviewed orders --- about ten a day on this
volume of 207 a day --- two listed customers in three went on to rate the order
1--2 stars, 6.1 times random, covering 30.7\% of such ratings; precision stays
above 0.125, the break-even only if every contact worked, at every capacity
(Figure~\ref{fig:capacity}a).
\textbf{Not everyone on that list can still be reached at delivery.} At ten a day
its 0.572 precision splits three ways: per 100 listed customers, 26.4 rated an
order that never arrived, 22.0 had already written while waiting for a late
parcel, and 8.8 could still be reached when the parcel came
(Appendix~\ref{app:timing}). Built as a team would actually work it --- each
morning, the previous day's deliveries whose customer has not yet rated the
order --- the daily top ten reaches 26 future 1--2 star reviewers in 100, 3.2
times the 8.1\% base rate of that pool; late and undelivered orders are reached
on their promised date instead (Section~\ref{sec:stakeholders}). \textbf{The
derived late flag dominates the direct classifier at every capacity}, by a factor
of 1.4 to 2.8 on precision, \textbf{and never clears the break-even for
an expensive action}: its precision peaks at 0.274 and falls as $K$ grows, under
the 0.300 a R\$12 expedite would need.

\subsection{Does the third family earn its place, and what is deployed}
\label{sec:results-ensemble}

Family C is selected on both problems, but whether it earns that differs. On
\textbf{lead time} it earns its place: the Voting blend improves on the better of
its members by 0.24 days of MAE (5.9\%) and is the only model in
Table~\ref{tab:reg} with a positive $R^2$, for the extrapolation reason above. On
\textbf{review risk} it does not: Stacking wins the rank vote by a nose (mean
rank 1.67 against 2.00, ahead on three of the six windows), then matches the
boosted trees to three decimals on ROC-AUC and comes in 0.001 \emph{below} them
on PR-AUC, although the ensemble was \emph{flattered}: stacking's inner folds
must partition the training set and so cannot be an expanding window
(Appendix~\ref{app:hyperparameters}).

\textbf{What is deployed.} The Phase 2 report named the deployed review model
inconsistently; the artifacts settle it. The served models, read back from the pickles
(Table~\ref{tab:deployed}), are the \textbf{VotingRegressor} for lead time, the
derived route with its two-parameter calibration for the late flag (not the
rule's pick; Section~\ref{sec:results-late}), and the \textbf{StackingClassifier}
for review risk, whose meta-learner weights the tree member 13.8 times the linear
one. The single boosted model would lose nothing measurable, but switching to it
after seeing test scores would be selection on test; a validation check ties
artifact, metadata and this text to the same estimator.

\subsection{Error analysis: where the models fail}
\label{sec:results-segments}

Across geography the review model is stable: ROC-AUC runs 0.744 to 0.778 across
the five regions, with lift over each region's floor between 3.6$\times$ and
5.0$\times$ (Figure~\ref{fig:segments}b). The regression is not: MAE runs from
3.42 days in the Southeast to \textbf{6.84 days in the North}, and the bias is
positive everywhere, $+1.15$ days in the Southeast to $+4.77$ in the North, in
every test month (Figures~\ref{fig:segments}a and \ref{fig:resid}): a model
fitted on a slower past over-predicts once fulfilment improves, the strongest
argument for refitting on a rolling basis. Table~\ref{tab:tvt}
(Appendix~\ref{app:tvt}) separates that drift from over-fitting: the blend's
in-sample MAE (4.72) sits level with validation (4.78) and above test (3.90), and
the stacked review model shows a tree ensemble's modest in-sample optimism
(ROC-AUC 0.836 against 0.762 and 0.749 on the periods it never saw).
